
\begingroup
\shorthandoff{<>}% Babel español trata de otro modo los delimitadores de curvatura de Xy-pic.
\begin{proposition}
\label{XIII.1.12}
Sean $S$ un esquema y $f\colon X \to Y$ un $S$-morfismo.
\begin{enumerate}
\item[1)] Sea
$$
\xymatrix@C=.5cm{ \Phi \ar[r] & \Phi_1 \ar@<2pt>[r] \ar@<-2pt>[r] &
\Phi_2 }
$$
un diagrama exacto de stacks sobre~$X$. Supongamos que $(\Phi_1, f)$ es
cohomológicamente propio relativamente a~$S$ en dimensión
$\leqslant 0$ y que $(\Phi_2, f)$ es cohomológicamente propio relativamente
a~$S$ en dimensión $\leqslant -1$. Entonces $(\Phi,f)$ es
cohomológicamente propio relativamente a~$S$ en dimensión $\leqslant 0$.

\item[2)] Sea
$$
\xymatrix@C=.5cm{ \Phi \ar[r] & \Phi_1 \ar@<2pt>[r] \ar@<-2pt>[r] &
\Phi_2 \ar@<4pt>[r] \ar[r] \ar@<-4pt>[r] & \Phi_3 }
$$
un diagrama exacto de stacks sobre~$X$. Supongamos que
\ifthenelse{\boolean{orig}}
{$(\Phi,f)$}
{$(\Phi_1,f)$}
es cohomológicamente propio relativamente a~$S$ en dimensión
$\leqslant 1$, que $(\Phi_2, f)$ es cohomológicamente propio
\marginpar{362}
relativamente a~$S$ en dimensión $\leqslant 0$ y que
$(\Phi_3, f)$ es cohomológicamente propio
\ifthenelse{\boolean{orig}}
{}
{relativamente}
a~$S$ en dimensión $\leqslant -1$. Entonces $(\Phi,f)$ es
cohomológicamente propio relativamente a~$S$ en dimensión $\leqslant 1$.
\end{enumerate}
\end{proposition}

Para todo $S$-esquema~$S'$, consideremos el diagrama conmutativo siguiente,
en el que todos los cuadrados son cartesianos:
$$
\xymatrix{ X' \ar[r]^{f'} \ar[d]_h & Y' \ar[r] \ar[d]_g & S' \ar[d] \\
X \ar[r]^f & Y \ar[r] & S }
$$

Demostremos 2), pues la demostración de 1) es análoga. Como los funtores de
imagen directa y de imagen inversa transforman
\ifthenelse{\boolean{orig}}
{diagrama exacto de stacks en diagrama exacto}
{un diagrama exacto de stacks en un diagrama exacto}
\eqref{XIII.1.10.4}, se tiene el morfismo siguiente
de diagramas exactos de stacks:
$$
\xymatrix{ g^* f_* \Phi \ar[r]^\pi \ar[d]_{\varphi} & g^*
f_* \Phi_1 \ar@<2pt>[r]^{\pi_1} \ar@<-2pt>[r]_{\pi_2}
\ar[d]_{\varphi_1} & g^* f_* \Phi_2 \ar@<4pt>[r] \ar[r]
\ar@<-4pt>[r] \ar[d]_{\varphi_2} & g^* f_* \Phi_3
\ar[d]_{\varphi_3} \\ f'_* h^* \Phi \ar[r] & f'_* h^*
\Phi_1 \ar@<2pt>[r] \ar@<-2pt>[r] & f'_* h^* \Phi_2
\ar@<4pt>[r] \ar[r] \ar@<-4pt>[r] & f'_* h^* \Phi_3 }
\text{.}
$$
Por hipótesis, $\varphi_1$ es una equivalencia de categorías,
$\varphi_2$ es plenamente fiel y $\varphi_3$ es fiel. Se deduce, por tanto,
del diagrama anterior que~$\varphi$ es una equivalencia.
\endgroup
